% ============================================================
% SECCIÓN A: Gestión de riesgos, sprint backlogs y cronograma
% ============================================================
\chapter{Sección A: Gestión del Proyecto}

% ------------------------------------------------------------
\section{Gestión de Riesgos}

\subsection{Marco de Identificación y Análisis}

Los riesgos del proyecto se identificaron mediante tres técnicas complementarias: listas de chequeo de riesgos típicos de proyectos de software, categorías documentadas de riesgo por módulo/actor del sistema y sesiones de tormenta de ideas del equipo, contrastadas contra los artefactos de planificación del proyecto (historias de usuario y cronograma PERT/CPM, ver Sección~\ref{sec:cronograma-a}).

Cada riesgo identificado se evalúa según dos dimensiones cualitativas de 5 niveles:

\begin{itemize}[leftmargin=*]
    \item \textbf{Probabilidad de ocurrencia:} (1) Raro [1--20\%], (2) Poco probable [20--40\%], (3) Moderado [40--60\%], (4) Probable [60--80\%], (5) Casi seguro [80--100\%].
    \item \textbf{Impacto potencial:} (1) Insignificante --- mínimo---, (2) Menor --- podría no afectar la satisfacción del cliente, presupuesto, cronograma o desempeño---, (3) Significativo --- potencial insatisfacción del cliente, sobrecosto, retraso o degradación del sistema---, (4) Mayor --- efecto significativo sobre cliente, cronograma, presupuesto o desempeño---, (5) Severo --- severa insatisfacción del cliente, sobrecosto grave, retraso grave o degradación severa del sistema.
\end{itemize}

El nivel de riesgo resultante (Probabilidad $\times$ Impacto) se clasifica de acuerdo con la matriz de calor 5$\times$5 de la Tabla~\ref{tab:matriz-riesgo}, la cual determina el orden de prioridad de atención: primero los riesgos de alta probabilidad y alto impacto, luego los de alto impacto con baja probabilidad, después los de bajo impacto con alta probabilidad, y por último los de bajo impacto y baja probabilidad.

\begin{table}[H]
\centering
\caption{Matriz de riesgo 5$\times$5 (Probabilidad $\times$ Impacto)}
\label{tab:matriz-riesgo}
\resizebox{\textwidth}{!}{%
\begin{tabular}{|c|c|c|c|c|c|}
\hline
\rowcolor{cajagris}
\textbf{Prob. $\downarrow$ / Impacto $\rightarrow$} & \textbf{Insign. (1)} & \textbf{Menor (2)} & \textbf{Signif. (3)} & \textbf{Mayor (4)} & \textbf{Severo (5)} \\
\hline
\textbf{5 -- Casi seguro} & \cellcolor{riskmedio}Medio 5 & \cellcolor{riskalto}Alto 10 & \cellcolor{riskmuyalto}Muy alto 15 & \cellcolor{riskextremo}Extremo 20 & \cellcolor{riskextremo}Extremo 25 \\
\hline
\textbf{4 -- Probable} & \cellcolor{riskmedio}Medio 4 & \cellcolor{riskmedio}Medio 8 & \cellcolor{riskalto}Alto 12 & \cellcolor{riskmuyalto}Muy alto 16 & \cellcolor{riskextremo}Extremo 20 \\
\hline
\textbf{3 -- Moderado} & \cellcolor{riskbajo}Bajo 3 & \cellcolor{riskmedio}Medio 6 & \cellcolor{riskmedio}Medio 9 & \cellcolor{riskalto}Alto 12 & \cellcolor{riskmuyalto}Muy alto 15 \\
\hline
\textbf{2 -- Poco probable} & \cellcolor{riskmuybajo}Muy bajo 2 & \cellcolor{riskbajo}Bajo 4 & \cellcolor{riskmedio}Medio 6 & \cellcolor{riskmedio}Medio 8 & \cellcolor{riskalto}Alto 10 \\
\hline
\textbf{1 -- Raro} & \cellcolor{riskmuybajo}Muy bajo 1 & \cellcolor{riskmuybajo}Muy bajo 2 & \cellcolor{riskbajo}Bajo 3 & \cellcolor{riskmedio}Medio 4 & \cellcolor{riskmedio}Medio 5 \\
\hline
\end{tabular}%
}
\end{table}

Para el manejo de cada riesgo se aplica una (o una combinación) de las siguientes seis estrategias:

\begin{enumerate}[leftmargin=*]
    \item \textbf{Evitar el riesgo:} remover por completo la causa del riesgo.
    \item \textbf{Reducción de la probabilidad:} disminuir la probabilidad de que el riesgo ocurra.
    \item \textbf{Reducción de impacto:} disminuir la severidad de las consecuencias si el riesgo ocurre.
    \item \textbf{Transferencia del riesgo:} trasladar el impacto (típicamente financiero) a un tercero, p. ej. mediante librerías/servicios ya auditados o subcontratación.
    \item \textbf{Planes de contingencia:} acciones predefinidas que se ejecutan si la amenaza se materializa.
    \item \textbf{Aceptación del riesgo:} riesgos de bajo impacto y baja probabilidad que se monitorean sin actuar activamente sobre ellos, salvo que su probabilidad o impacto cambie.
\end{enumerate}

\subsection{Registro de Riesgos del Proyecto}

Los riesgos se identificaron y organizaron por módulo/actor del sistema (Gestor de Bodega, Cliente, Administrador, Organización y Leodeguita), cada uno vinculado a la historia de usuario que lo origina. Se identificaron 32 riesgos funcionales; ninguno alcanza los niveles ``Muy alto'' o ``Extremo'': 4 son ``Alto'', 25 ``Medio'' y 3 ``Bajo'' (Tabla~\ref{tab:riesgos-resumen}). Para cada riesgo se documenta, además de su evaluación, cómo el sistema lo previene o mitiga en la práctica (Tabla~\ref{tab:riesgos-plan}).

\begin{table}[H]
\centering
\caption{Resumen de riesgos por nivel}
\label{tab:riesgos-resumen}
\begin{tabular}{|l|c|}
\hline
\rowcolor{cajagris}
\textbf{Nivel} & \textbf{Cantidad} \\
\hline
\cellcolor{riskmuybajo}Muy bajo & 0 \\ \hline
\cellcolor{riskbajo}Bajo & 3 \\ \hline
\cellcolor{riskmedio}Medio & 25 \\ \hline
\cellcolor{riskalto}Alto & 4 \\ \hline
\cellcolor{riskmuyalto}Muy alto & 0 \\ \hline
\cellcolor{riskextremo}Extremo & 0 \\ \hline
\textbf{Total} & \textbf{32} \\ \hline
\end{tabular}
\end{table}

{\small\setlength{\tabcolsep}{4pt}
\begin{longtable}{|M{1.2cm}|M{2.5cm}|L{5.0cm}|M{2.0cm}|M{2.0cm}|M{1.5cm}|}
\caption{Identificación y evaluación de riesgos} \label{tab:riesgos-eval} \\
\hline
\rowcolor{cajagris}
\textbf{ID} & \textbf{Módulo} & \textbf{Riesgo} & \textbf{Prob.} & \textbf{Impacto} & \textbf{Nivel} \\
\hline
\endfirsthead
\multicolumn{6}{c}{\textit{Tabla~\ref{tab:riesgos-eval} (continuación)}} \\
\hline
\rowcolor{cajagris}
\textbf{ID} & \textbf{Módulo} & \textbf{Riesgo} & \textbf{Prob.} & \textbf{Impacto} & \textbf{Nivel} \\
\hline
\endhead
\hline \multicolumn{6}{r}{\textit{continúa en la siguiente página}} \\
\endfoot
\hline
\endlastfoot

RG-01 & Gestor de Bodega & Permiso de bomberos vencido o inválido al registrar la bodega (HUG-04). & 3 -- Moderado & 4 -- Mayor & \cellcolor{riskalto}Alto (12) \\
\hline
RG-02 & Gestor de Bodega & Registro de bodega con datos incompletos o inválidos (HUG-04). & 4 -- Probable & 2 -- Menor & \cellcolor{riskmedio}Medio (8) \\
\hline
RG-03 & Gestor de Bodega & Cancelación de reserva ya iniciada o finalizada (HUG-06). & 2 -- Poco probable & 3 -- Significativo & \cellcolor{riskmedio}Medio (6) \\
\hline
RG-04 & Gestor de Bodega & Eliminación de bodega con reservas activas o futuras (HUG-07). & 2 -- Poco probable & 4 -- Mayor & \cellcolor{riskmedio}Medio (8) \\
\hline
RG-05 & Gestor de Bodega & Edición de bodega afecta reservas confirmadas (HUG-08). & 3 -- Moderado & 3 -- Significativo & \cellcolor{riskmedio}Medio (9) \\
\hline
RG-06 & Gestor de Bodega & No recibir notificación de pago confirmado (HUG-02). & 2 -- Poco probable & 3 -- Significativo & \cellcolor{riskmedio}Medio (6) \\
\hline
RG-07 & Gestor de Bodega & Chat no disponible tras finalizar el alquiler (HUG-03). & 5 -- Casi seguro & 1 -- Insignificante & \cellcolor{riskmedio}Medio (5) \\
\hline
RC-01 & Cliente & Reserva con fechas solapadas con otra confirmada (HUC-03). & 3 -- Moderado & 3 -- Significativo & \cellcolor{riskmedio}Medio (9) \\
\hline
RC-02 & Cliente & No completar el pago dentro del tiempo de hold (HUC-05). & 3 -- Moderado & 3 -- Significativo & \cellcolor{riskmedio}Medio (9) \\
\hline
RC-03 & Cliente & Endpoint de inventario externo no responde o es inválido (HUC-10, HUC-12). & 4 -- Probable & 3 -- Significativo & \cellcolor{riskalto}Alto (12) \\
\hline
RC-04 & Cliente & Mapeo de campos de inventario incorrecto (HUC-11). & 3 -- Moderado & 3 -- Significativo & \cellcolor{riskmedio}Medio (9) \\
\hline
RC-05 & Cliente & Importación de Excel con filas incompletas o formato incorrecto (HUC-13). & 4 -- Probable & 2 -- Menor & \cellcolor{riskmedio}Medio (8) \\
\hline
RC-06 & Cliente & Intento de cancelar reserva ya iniciada (HUC-07). & 3 -- Moderado & 3 -- Significativo & \cellcolor{riskmedio}Medio (9) \\
\hline
RC-07 & Cliente & Invitación de organización expirada (HUC-14). & 3 -- Moderado & 2 -- Menor & \cellcolor{riskmedio}Medio (6) \\
\hline
RC-08 & Cliente & Reserva o consulta sin sesión activa (HUC-03, HUC-04). & 3 -- Moderado & 1 -- Insignificante & \cellcolor{riskbajo}Bajo (3) \\
\hline
RA-01 & Administrador & Aprobar bodega con permiso de bomberos inválido (HUA-01). & 2 -- Poco probable & 5 -- Severo & \cellcolor{riskalto}Alto (10) \\
\hline
RA-02 & Administrador & Bloquear cuenta de usuario legítimo por error (HUA-03). & 2 -- Poco probable & 4 -- Mayor & \cellcolor{riskmedio}Medio (8) \\
\hline
RA-03 & Administrador & Métricas del dashboard inexactas o desactualizadas (HUA-02). & 3 -- Moderado & 3 -- Significativo & \cellcolor{riskmedio}Medio (9) \\
\hline
RA-04 & Administrador & Resolver reporte sin respuesta al cliente (HUA-04). & 3 -- Moderado & 3 -- Significativo & \cellcolor{riskmedio}Medio (9) \\
\hline
RA-05 & Administrador & Acumulación de reportes sin resolver (HUA-02, HUA-04). & 3 -- Moderado & 3 -- Significativo & \cellcolor{riskmedio}Medio (9) \\
\hline
RO-01 & Organización & Crear organización con RUC duplicado (HUE-04). & 2 -- Poco probable & 3 -- Significativo & \cellcolor{riskmedio}Medio (6) \\
\hline
RO-02 & Organización & Invitar a usuarios con rol gestor en vez de cliente (HUE-01). & 3 -- Moderado & 2 -- Menor & \cellcolor{riskmedio}Medio (6) \\
\hline
RO-03 & Organización & Miembro elimina reservas de la organización sin permiso (HUE-06). & 2 -- Poco probable & 4 -- Mayor & \cellcolor{riskmedio}Medio (8) \\
\hline
RO-04 & Organización & Reserva empresarial hecha en contexto personal por error (HUE-05). & 4 -- Probable & 3 -- Significativo & \cellcolor{riskalto}Alto (12) \\
\hline
RO-05 & Organización & Cambio de contexto erróneo personal/organización (HUE-02). & 3 -- Moderado & 2 -- Menor & \cellcolor{riskmedio}Medio (6) \\
\hline
RO-06 & Organización & Miembro eliminado pierde acceso a la organización (HUE-03). & 3 -- Moderado & 2 -- Menor & \cellcolor{riskmedio}Medio (6) \\
\hline
RL-01 & Leodeguita & Inicio de sesión con credenciales incorrectas (HUL-01). & 4 -- Probable & 1 -- Insignificante & \cellcolor{riskmedio}Medio (4) \\
\hline
RL-02 & Leodeguita & Publicar bodega con nombre duplicado (HUL-03). & 3 -- Moderado & 2 -- Menor & \cellcolor{riskmedio}Medio (6) \\
\hline
RL-03 & Leodeguita & Campos incompletos en publicación desde el móvil (HUL-03). & 4 -- Probable & 2 -- Menor & \cellcolor{riskmedio}Medio (8) \\
\hline
RL-04 & Leodeguita & Espacio no disponible mientras se visualiza (HUL-06). & 3 -- Moderado & 2 -- Menor & \cellcolor{riskmedio}Medio (6) \\
\hline
RL-05 & Leodeguita & Cliente sin compañías asociadas en Leodeguita (HUL-02). & 3 -- Moderado & 1 -- Insignificante & \cellcolor{riskbajo}Bajo (3) \\
\hline
RL-06 & Leodeguita & Listado de bodegas o espacios vacío en el móvil (HUL-04, HUL-05). & 3 -- Moderado & 1 -- Insignificante & \cellcolor{riskbajo}Bajo (3) \\
\hline
\end{longtable}
}

{\setlength{\tabcolsep}{4pt}
\begin{longtable}{|M{1.1cm}|M{2.4cm}|L{11.6cm}|}
\caption{Plan de respuesta a riesgos} \label{tab:riesgos-plan} \\
\hline
\rowcolor{cajagris}
\textbf{ID} & \textbf{Estrategia} & \textbf{Aplicación en el sistema} \\
\hline
\endfirsthead
\multicolumn{3}{c}{\textit{Tabla~\ref{tab:riesgos-plan} (continuación)}} \\
\hline
\rowcolor{cajagris}
\textbf{ID} & \textbf{Estrategia} & \textbf{Aplicación en el sistema} \\
\hline
\endhead
\hline \multicolumn{3}{r}{\textit{continúa en la siguiente página}} \\
\endfoot
\hline
\endlastfoot

RG-01 & Reducción de probabilidad & Validar la vigencia del permiso antes de enviar; el sistema exige adjuntar el documento obligatoriamente y lo envía a revisión del administrador (HUA-01). \\
\hline
RG-02 & Reducción de probabilidad & Validación de campos obligatorios en el formulario con mensajes de error claros antes de permitir el envío. \\
\hline
RG-03 & Evitar & El sistema bloquea la cancelación de reservas en curso o finalizadas y muestra un mensaje de error indicando que no es posible. \\
\hline
RG-04 & Evitar & El sistema bloquea la eliminación cuando existen reservas activas o futuras y solicita gestionarlas antes de eliminar. \\
\hline
RG-05 & Reducción de impacto & Los cambios de precio u otros datos aplican solo a nuevas reservas; las existentes conservan las condiciones originales acordadas. \\
\hline
RG-06 & Planes de contingencia & El gestor puede verificar contratos activos en su panel; el sistema reenvía la notificación y marca el contrato como activo. \\
\hline
RG-07 & Aceptación & El sistema deshabilita el chat en modo solo lectura y conserva el historial para referencia; no se requiere acción adicional. \\
\hline
RC-01 & Evitar & El sistema detecta conflictos de fechas, bloquea la reserva y resalta los periodos ocupados en el calendario de disponibilidad. \\
\hline
RC-02 & Planes de contingencia & El sistema libera la disponibilidad retenida, no finaliza la reserva y permite al cliente iniciar una nueva cuando lo desee. \\
\hline
RC-03 & Planes de contingencia & El sistema muestra los últimos datos disponibles con su fecha de consulta y ofrece importar desde Excel como alternativa. \\
\hline
RC-04 & Reducción de probabilidad & El sistema valida los campos obligatorios (nombre y cantidad) y muestra una vista previa del mapeo antes de guardar. \\
\hline
RC-05 & Reducción de impacto & El sistema importa las filas válidas, omite las incorrectas y muestra un reporte detallado de las filas ignoradas y el motivo. \\
\hline
RC-06 & Transferencia & El sistema bloquea la cancelación y deriva al cliente al administrador para resolver incidencias en curso. \\
\hline
RC-07 & Planes de contingencia & El sistema archiva la invitación y sugiere al cliente solicitar una nueva al administrador de la organización. \\
\hline
RC-08 & Evitar & El sistema redirige al login y, tras autenticar, retorna automáticamente a la bodega o reserva seleccionada. \\
\hline
RA-01 & Reducción de impacto & El administrador revisa el documento; el rechazo exige razón obligatoria, se notifica al gestor y debe subir un permiso válido para reiniciar. \\
\hline
RA-02 & Reducción de impacto & El bloqueo requiere razón obligatoria y se registra en bitácora de auditoría; la cuenta puede reactivarse posteriormente. \\
\hline
RA-03 & Reducción de probabilidad & El dashboard consulta columnas indexadas y filtra métricas por periodo con gráficos comparativos mensuales. \\
\hline
RA-04 & Evitar & El sistema bloquea la resolución si el campo de respuesta está vacío y solicita el mensaje antes de cerrar el reporte. \\
\hline
RA-05 & Planes de contingencia & El dashboard muestra un indicador numérico de reportes pendientes clasificados por prioridad con enlace directo al módulo. \\
\hline
RO-01 & Evitar & El sistema valida la unicidad del RUC (13 dígitos) y bloquea el registro si ya existe otra organización con ese RUC. \\
\hline
RO-02 & Evitar & El sistema bloquea la invitación a gestores de bodega y muestra que solo se pueden invitar usuarios con rol cliente. \\
\hline
RO-03 & Evitar & El botón de eliminar se oculta a miembros; los intentos por URL directa devuelven error 403 y se registran en auditoría. \\
\hline
RO-04 & Reducción de probabilidad & El selector de contexto indica claramente el contexto activo; la reserva se vincula al contexto vigente al confirmar. \\
\hline
RO-05 & Reducción de impacto & El indicador visual muestra el contexto activo y el panel se recarga con los datos específicos de ese contexto. \\
\hline
RO-06 & Aceptación & El miembro es desvinculado y notificado; el historial de reservas que realizó se preserva para auditoría. \\
\hline
RL-01 & Reducción de probabilidad & El sistema valida los campos, resalta los obligatorios vacíos y muestra un mensaje de error claro sin permitir el acceso. \\
\hline
RL-02 & Evitar & El sistema advierte que ya existe una bodega con ese nombre del mismo gestor y bloquea el registro. \\
\hline
RL-03 & Reducción de probabilidad & El sistema resalta los campos faltantes con un mensaje descriptivo y no realiza el registro hasta completarlos. \\
\hline
RL-04 & Planes de contingencia & El sistema muestra un aviso de que el espacio ya no está disponible y lo retira de la lista activa al refrescar. \\
\hline
RL-05 & Aceptación & El sistema redirige a la vista del rol cliente y no muestra el selector de compañía al no tener organizaciones asociadas. \\
\hline
RL-06 & Aceptación & El sistema muestra un mensaje indicando que no hay bodegas/espacios y ofrece la opción de publicar o explorar. \\
\hline
\end{longtable}
}

\clearpage

% ------------------------------------------------------------
\section{Sprint Backlogs}
El proyecto se organizó en \textbf{5 sprints}, con una duración total planificada de 52 días según el cronograma PERT/CPM (Sección~\ref{sec:cronograma-a}), gestionados mediante un tablero ágil tipo Scrum en \textit{Jira}. Cada historia de usuario se prioriza con la escala MoSCoW (Must/Should/Could Have) y se estima en puntos de historia.

\subsection{Sprint 1 --- Fundamentos de Autenticación y Gestión de Almacenes}
\textbf{Objetivo del Sprint:} Establecer el inicio de sesión compartido web/móvil y las operaciones básicas de registro, edición, eliminación y moderación de almacenes.\\
\textbf{Duración:} Día 1 -- Día 11 (11 días)

{\setlength{\tabcolsep}{4pt}
\begin{longtable}{|M{1.6cm}|L{7.0cm}|M{1.4cm}|M{2.4cm}|M{1.8cm}|}
\caption{Sprint Backlog -- Sprint 1} \label{tab:sprint1} \\
\hline
\rowcolor{cajagris}
\textbf{ID} & \textbf{Descripción (Historia de Usuario)} & \textbf{Puntos} & \textbf{Responsable} & \textbf{Estado} \\
\hline
\endfirsthead
\multicolumn{5}{c}{\textit{Tabla~\ref{tab:sprint1} (continuación)}} \\
\hline
\rowcolor{cajagris}
\textbf{ID} & \textbf{Descripción (Historia de Usuario)} & \textbf{Puntos} & \textbf{Responsable} & \textbf{Estado} \\
\hline
\endhead
\hline \multicolumn{5}{r}{\textit{continúa en la siguiente página}} \\
\endfoot
\hline
\endlastfoot

HUG-04 & Como Gestor de Almacenamiento, quiero registrar un nuevo almacén (dimensiones, ubicación, fotos, tarifa y permiso de bomberos vigente), para ofrecer mi espacio en alquiler. & 5 & Equipo Leodega & Pendiente \\ \hline
HUG-07 & Como Gestor, quiero eliminar un almacén sin reservas activas o futuras, para retirar espacios que ya no ofrezco. & 3 & Equipo Leodega & Pendiente \\ \hline
HUG-08 & Como Gestor, quiero editar la información de un almacén ya publicado, para mantener actualizada mi oferta sin recrearlo. & 5 & Equipo Leodega & Pendiente \\ \hline
HUA-01 & Como Administrador, quiero revisar y validar el permiso de bomberos subido por el Gestor, para asegurar el cumplimiento legal antes de habilitar el espacio. & 5 & Equipo Leodega & Pendiente \\ \hline
HUA-03 & Como Administrador, quiero bloquear o reactivar cuentas de usuario, para actuar como moderador ante incumplimientos de las políticas de uso. & 5 & Equipo Leodega & Pendiente \\ \hline
HUL-01 & Como usuario (Administrador, Gestor o Cliente), quiero iniciar sesión en Leodeguita con mis credenciales de Leodega, para acceder a mis datos sin crear una cuenta aparte. & 3 & Equipo Leodega & Pendiente \\ \hline
HUL-03 & Como Gestor, quiero registrar un almacén disponible con el mínimo de pasos, para publicarlo rápidamente desde mi dispositivo móvil. & 5 & Equipo Leodega & Pendiente \\ \hline
HUL-04 & Como Gestor, quiero ver la lista de almacenes que he publicado, para monitorear mis espacios y su estado actual. & 2 & Equipo Leodega & Pendiente \\ \hline
\end{longtable}
}
\clearpage

\subsection{Sprint 2 --- Catálogo, Reservas y Pago}
\textbf{Objetivo del Sprint:} Habilitar la búsqueda y reserva de almacenes por parte del cliente, el pago simulado y la gestión de reservas por parte del Gestor.\\
\textbf{Duración:} Día 12 -- Día 21 (10 días)

{\setlength{\tabcolsep}{4pt}
\begin{longtable}{|M{1.6cm}|L{7.0cm}|M{1.4cm}|M{2.4cm}|M{1.8cm}|}
\caption{Sprint Backlog -- Sprint 2} \label{tab:sprint2} \\
\hline
\rowcolor{cajagris}
\textbf{ID} & \textbf{Descripción (Historia de Usuario)} & \textbf{Puntos} & \textbf{Responsable} & \textbf{Estado} \\
\hline
\endfirsthead
\multicolumn{5}{c}{\textit{Tabla~\ref{tab:sprint2} (continuación)}} \\
\hline
\rowcolor{cajagris}
\textbf{ID} & \textbf{Descripción (Historia de Usuario)} & \textbf{Puntos} & \textbf{Responsable} & \textbf{Estado} \\
\hline
\endhead
\hline \multicolumn{5}{r}{\textit{continúa en la siguiente página}} \\
\endfoot
\hline
\endlastfoot

HUG-02 & Como Gestor, quiero recibir una notificación cuando un cliente confirme el pago de una reserva, para saber en tiempo real qué contratos se formalizaron. & 3 & Equipo Leodega & Pendiente \\ \hline
HUG-05 & Como Gestor, quiero ver todas las reservas de mis almacenes con fecha, cliente y estado de pago, para gestionar el uso de mi espacio. & 3 & Equipo Leodega & Pendiente \\ \hline
HUG-06 & Como Gestor, quiero cancelar una reserva pagada indicando un motivo obligatorio, para liberar mi espacio ante imprevistos manteniendo transparencia con el cliente. & 3 & Equipo Leodega & Pendiente \\ \hline
HUC-01 & Como Cliente, quiero buscar y filtrar almacenes disponibles por ubicación, tamaño y precio, para encontrar el espacio que mejor se ajuste a mis necesidades y presupuesto. & 5 & Equipo Leodega & Pendiente \\ \hline
HUC-02 & Como Cliente, quiero ver el detalle completo de un almacén antes de reservar, para evaluar si cumple mis requisitos de condiciones, capacidad, precio y ubicación. & 3 & Equipo Leodega & Pendiente \\ \hline
HUC-03 & Como Cliente, quiero reservar un almacén disponible seleccionando el período de alquiler deseado, para asegurar el espacio que necesito en las fechas requeridas. & 5 & Equipo Leodega & Pendiente \\ \hline
HUC-05 & Como Cliente, quiero completar un proceso de pago simulado para confirmar y activar mi reserva, para formalizar el acuerdo de alquiler. & 5 & Equipo Leodega & Pendiente \\ \hline
HUL-06 & Como Cliente, quiero ver la información detallada de un espacio disponible desde el móvil, para evaluar si cumple mis necesidades antes de alquilarlo. & 2 & Equipo Leodega & Pendiente \\ \hline
\end{longtable}
}
\clearpage

\subsection{Sprint 3 --- Organizaciones (Multiempresa)}
\textbf{Objetivo del Sprint:} Incorporar la operación multiempresa: creación de organizaciones, invitación de miembros, cambio de contexto y reservas a nombre de la organización.\\
\textbf{Duración:} Día 22 -- Día 32 (11 días)

{\setlength{\tabcolsep}{4pt}
\begin{longtable}{|M{1.6cm}|L{7.0cm}|M{1.4cm}|M{2.4cm}|M{1.8cm}|}
\caption{Sprint Backlog -- Sprint 3} \label{tab:sprint3} \\
\hline
\rowcolor{cajagris}
\textbf{ID} & \textbf{Descripción (Historia de Usuario)} & \textbf{Puntos} & \textbf{Responsable} & \textbf{Estado} \\
\hline
\endfirsthead
\multicolumn{5}{c}{\textit{Tabla~\ref{tab:sprint3} (continuación)}} \\
\hline
\rowcolor{cajagris}
\textbf{ID} & \textbf{Descripción (Historia de Usuario)} & \textbf{Puntos} & \textbf{Responsable} & \textbf{Estado} \\
\hline
\endhead
\hline \multicolumn{5}{r}{\textit{continúa en la siguiente página}} \\
\endfoot
\hline
\endlastfoot

HUC-14 & Como Cliente, quiero aceptar o rechazar una invitación de una organización para unirme como miembro, para decidir si opero bajo su nombre en la plataforma. & 3 & Equipo Leodega & Pendiente \\ \hline
HUE-01 & Como Administrador de Organización, quiero invitar a usuarios clientes a unirse a mi organización como miembros, para que mis colaboradores operen bajo el nombre de la empresa. & 5 & Equipo Leodega & Pendiente \\ \hline
HUE-02 & Como Cliente perteneciente a una organización, quiero cambiar entre mi contexto personal y el de cada organización desde un selector, para operar bajo el nombre correcto sin cerrar sesión. & 5 & Equipo Leodega & Pendiente \\ \hline
HUE-03 & Como Administrador de Organización, quiero acceder a un panel administrativo con miembros, reservas activas e historial, para tener visibilidad y control total de las operaciones de mi empresa. & 5 & Equipo Leodega & Pendiente \\ \hline
HUE-04 & Como Administrador de Organización, quiero crear una organización desde mi cuenta de cliente ingresando sus datos (razón social, RUC y correo), para gestionar alquileres a nombre de mi organización. & 5 & Equipo Leodega & Pendiente \\ \hline
HUE-05 & Como Cliente perteneciente a una organización, quiero reservar un almacén disponible a nombre de la organización, para asegurar espacio de almacenamiento para la empresa. & 5 & Equipo Leodega & Pendiente \\ \hline
HUE-06 & Como Cliente perteneciente a una organización, quiero que el sistema me impida eliminar almacenes reservados por la organización, para proteger los contratos empresariales de eliminaciones no autorizadas. & 3 & Equipo Leodega & Pendiente \\ \hline
HUL-02 & Como Cliente, quiero seleccionar la empresa bajo la cual operar dentro de Leodeguita, para ver solo los almacenes e ítems correspondientes a esa empresa. & 3 & Equipo Leodega & Pendiente \\ \hline
\end{longtable}
}
\clearpage

\subsection{Sprint 4 --- Mensajería, Reportes y Reseñas}
\textbf{Objetivo del Sprint:} Habilitar la comunicación cliente--gestor (consultas y chat de reserva activa), el reporte de incidentes y el sistema de calificaciones.\\
\textbf{Duración:} Día 33 -- Día 42 (10 días)

{\setlength{\tabcolsep}{4pt}
\begin{longtable}{|M{1.6cm}|L{7.0cm}|M{1.4cm}|M{2.4cm}|M{1.8cm}|}
\caption{Sprint Backlog -- Sprint 4} \label{tab:sprint4} \\
\hline
\rowcolor{cajagris}
\textbf{ID} & \textbf{Descripción (Historia de Usuario)} & \textbf{Puntos} & \textbf{Responsable} & \textbf{Estado} \\
\hline
\endfirsthead
\multicolumn{5}{c}{\textit{Tabla~\ref{tab:sprint4} (continuación)}} \\
\hline
\rowcolor{cajagris}
\textbf{ID} & \textbf{Descripción (Historia de Usuario)} & \textbf{Puntos} & \textbf{Responsable} & \textbf{Estado} \\
\hline
\endhead
\hline \multicolumn{5}{r}{\textit{continúa en la siguiente página}} \\
\endfoot
\hline
\endlastfoot

HUG-01 & Como Gestor, quiero recibir y responder mensajes de consulta previa de clientes interesados, para brindar información oportuna y generar confianza antes de la reserva formal. & 5 & Equipo Leodega & Pendiente \\ \hline
HUG-03 & Como Gestor, quiero tener un canal de chat activo con mis clientes durante el período de alquiler vigente, para coordinar accesos y resolver incidencias sin salir de la plataforma. & 5 & Equipo Leodega & Pendiente \\ \hline
HUC-04 & Como Cliente, quiero enviar un mensaje al Gestor de un almacén antes de reservar para aclarar dudas, para obtener información que no está en la descripción. & 3 & Equipo Leodega & Pendiente \\ \hline
HUC-06 & Como Cliente, quiero comunicarme con el Gestor mediante un chat vinculado a mi reserva activa, para coordinar accesos y resolver situaciones operativas menores. & 5 & Equipo Leodega & Pendiente \\ \hline
HUC-07 & Como Cliente, quiero ver y gestionar mis reservas activas e históricas desde mi panel personal, para tener visibilidad y control sobre mis contratos de alquiler vigentes. & 3 & Equipo Leodega & Pendiente \\ \hline
HUC-08 & Como Cliente, quiero reportar un problema relacionado con un almacén o una reserva activa desde mi panel personal, para notificar al Administrador incumplimientos, daños o situaciones anómalas. & 3 & Equipo Leodega & Pendiente \\ \hline
HUC-09 & Como Cliente, quiero calificar y dejar una reseña sobre el Gestor y el almacén al finalizar mi alquiler, para compartir mi experiencia y aportar a la reputación del Gestor. & 3 & Equipo Leodega & Pendiente \\ \hline
HUA-04 & Como Administrador, quiero revisar y dar seguimiento a los reportes que envían los clientes, pudiendo responder y cambiar su estado, para gestionar incidencias y mantener la calidad del servicio. & 3 & Equipo Leodega & Pendiente \\ \hline
\end{longtable}
}
\clearpage

\subsection{Sprint 5 --- Integración de Inventario y Métricas}
\textbf{Objetivo del Sprint:} Incorporar la integración con sistemas de inventario externos (conexión, mapeo y visualización), la importación vía Excel y el panel de métricas del administrador.\\
\textbf{Duración:} Día 43 -- Día 52 (10 días)

{\setlength{\tabcolsep}{4pt}
\begin{longtable}{|M{1.6cm}|L{7.0cm}|M{1.4cm}|M{2.4cm}|M{1.8cm}|}
\caption{Sprint Backlog -- Sprint 5} \label{tab:sprint5} \\
\hline
\rowcolor{cajagris}
\textbf{ID} & \textbf{Descripción (Historia de Usuario)} & \textbf{Puntos} & \textbf{Responsable} & \textbf{Estado} \\
\hline
\endfirsthead
\multicolumn{5}{c}{\textit{Tabla~\ref{tab:sprint5} (continuación)}} \\
\hline
\rowcolor{cajagris}
\textbf{ID} & \textbf{Descripción (Historia de Usuario)} & \textbf{Puntos} & \textbf{Responsable} & \textbf{Estado} \\
\hline
\endhead
\hline \multicolumn{5}{r}{\textit{continúa en la siguiente página}} \\
\endfoot
\hline
\endlastfoot

HUC-10 & Como Cliente, quiero registrar la URL del endpoint de mi sistema de inventario externo junto con la credencial de acceso, para que Leodega consulte automáticamente los productos almacenados sin modificar mi sistema. & 8 & Equipo Leodega & Pendiente \\ \hline
HUC-11 & Como Cliente, quiero configurar la correspondencia entre los campos devueltos por el endpoint de mi sistema externo y los campos de Leodega, para que mis datos de inventario se registren correctamente. & 8 & Equipo Leodega & Pendiente \\ \hline
HUC-12 & Como Cliente, quiero ver la lista actualizada de productos almacenados en mi espacio reservado, obtenida directamente de mi sistema externo, para controlar mi inventario sin depender de canales externos. & 5 & Equipo Leodega & Pendiente \\ \hline
HUC-13 & Como Cliente, quiero subir un archivo Excel con la lista de mis productos directamente a la plataforma, para registrar mi inventario rápidamente cuando mi sistema externo no expone un endpoint. & 5 & Equipo Leodega & Pendiente \\ \hline
HUA-02 & Como Administrador, quiero ver un panel de control con métricas clave del estado de la plataforma, para tomar decisiones informadas sobre la operación del negocio y monitorear la rentabilidad mensual. & 8 & Equipo Leodega & Pendiente \\ \hline
HUL-05 & Como Cliente, quiero ver todos los espacios de almacenamiento disponibles en una sola lista desde el móvil, para encontrar un espacio que se ajuste a mis necesidades. & 5 & Equipo Leodega & Pendiente \\ \hline
\end{longtable}
}
\clearpage

% ------------------------------------------------------------
\section{Cronograma del Proyecto}
\label{sec:cronograma-a}

El cronograma se construyó a partir de 38 actividades (\textbf{HST-001} a \textbf{HST-038}), derivadas de las historias de usuario de la Sección~2.2\footnote{Los identificadores HST-XXX numeran secuencialmente las actividades del análisis PERT/CPM y del diagrama Activity-on-Arrow; no coinciden con los códigos HUG/HUC/HUA/HUE/HUL de las historias de usuario porque responden a una convención de numeración independiente, propia del análisis de cronograma.}, con sus dependencias, duraciones en días y personal asignado. Para cada actividad se calculó el tiempo más próximo del evento (\textbf{EET}), la holgura total y el tiempo más lejano del evento (\textbf{LET}); las actividades con holgura cero pertenecen a la ruta crítica (\textbf{RC}) y se marcan con $\bullet$.

\subsection{Tabla de Actividades}
{\setlength{\tabcolsep}{4pt}
\begin{longtable}{|M{1.9cm}|M{1.3cm}|M{1.6cm}|M{1.3cm}|M{1.7cm}|M{1.1cm}|M{1.7cm}|M{1.1cm}|M{1.0cm}|}
\caption{Actividades, duraciones, dependencias y análisis CPM del proyecto} \label{tab:cronograma_tabla} \\
\hline
\rowcolor{cajagris}
\textbf{ID} & \textbf{Sprint} & \textbf{Predec.} & \textbf{Dur.} & \textbf{Personal} & \textbf{EET} & \textbf{Holgura} & \textbf{LET} & \textbf{RC} \\
\hline
\endfirsthead
\multicolumn{9}{c}{\textit{Tabla~\ref{tab:cronograma_tabla} (continuación)}} \\
\hline
\rowcolor{cajagris}
\textbf{ID} & \textbf{Sprint} & \textbf{Predec.} & \textbf{Dur.} & \textbf{Personal} & \textbf{EET} & \textbf{Holgura} & \textbf{LET} & \textbf{RC} \\
\hline
\endhead
\hline \multicolumn{9}{r}{\textit{continúa en la siguiente página}} \\
\endfoot
\hline
\endlastfoot

HST-001 & 1 & -- & 3 & 1 & 0 & 0 & 3 & $\bullet$ \\ \hline
HST-002 & 1 & HST-001 & 5 & 2 & 3 & 0 & 8 & $\bullet$ \\ \hline
HST-003 & 1 & HST-001 & 5 & 2 & 3 & 0 & 8 & $\bullet$ \\ \hline
HST-004 & 1 & HST-002 & 5 & 2 & 8 & 39 & 52 & \\ \hline
HST-005 & 1 & HST-002 & 3 & 1 & 8 & 41 & 52 & \\ \hline
HST-006 & 1 & HST-002, HST-003 & 2 & 1 & 8 & 42 & 52 & \\ \hline
HST-007 & 1 & HST-002, HST-003 & 5 & 2 & 8 & 0 & 13 & $\bullet$ \\ \hline
HST-008 & 1 & HST-001 & 5 & 2 & 3 & 36 & 44 & \\ \hline
HST-009 & 2 & HST-007 & 5 & 2 & 13 & 0 & 18 & $\bullet$ \\ \hline
HST-010 & 2 & HST-007 & 3 & 1 & 13 & 34 & 50 & \\ \hline
HST-011 & 2 & HST-009 & 3 & 1 & 18 & 0 & 21 & $\bullet$ \\ \hline
HST-012 & 2 & HST-010 & 2 & 1 & 16 & 34 & 52 & \\ \hline
HST-013 & 2 & HST-011 & 5 & 2 & 21 & 0 & 26 & $\bullet$ \\ \hline
HST-014 & 2 & HST-013 & 5 & 2 & 26 & 0 & 31 & $\bullet$ \\ \hline
HST-015 & 2 & HST-014 & 3 & 1 & 31 & 18 & 52 & \\ \hline
HST-016 & 2 & HST-014 & 3 & 1 & 31 & 15 & 49 & \\ \hline
HST-017 & 2 & HST-016 & 3 & 1 & 34 & 15 & 52 & \\ \hline
HST-018 & 3 & HST-001 & 5 & 2 & 3 & 28 & 36 & \\ \hline
HST-019 & 3 & HST-018 & 5 & 2 & 8 & 31 & 44 & \\ \hline
HST-020 & 3 & HST-018 & 3 & 1 & 8 & 41 & 52 & \\ \hline
HST-021 & 3 & HST-018 & 5 & 2 & 8 & 28 & 41 & \\ \hline
HST-022 & 3 & HST-021 & 3 & 1 & 13 & 28 & 44 & \\ \hline
HST-023 & 3 & HST-019, HST-022 & 5 & 2 & 16 & 23 & 49 & \\ \hline
HST-024 & 3 & HST-019, HST-014 & 5 & 2 & 31 & 13 & 49 & \\ \hline
HST-025 & 3 & HST-023, HST-024 & 3 & 1 & 36 & 13 & 52 & \\ \hline
HST-026 & 4 & HST-011 & 3 & 1 & 21 & 23 & 47 & \\ \hline
HST-027 & 4 & HST-026 & 5 & 2 & 24 & 23 & 52 & \\ \hline
HST-028 & 4 & HST-014 & 5 & 2 & 31 & 11 & 47 & \\ \hline
HST-029 & 4 & HST-028 & 5 & 2 & 36 & 11 & 52 & \\ \hline
HST-030 & 4 & HST-014 & 3 & 1 & 31 & 18 & 52 & \\ \hline
HST-031 & 4 & HST-014 & 3 & 1 & 31 & 15 & 49 & \\ \hline
HST-032 & 4 & HST-031 & 3 & 1 & 34 & 15 & 52 & \\ \hline
HST-033 & 4 & HST-014 & 3 & 1 & 31 & 18 & 52 & \\ \hline
HST-034 & 5 & HST-014 & 8 & 3 & 31 & 0 & 39 & $\bullet$ \\ \hline
HST-035 & 5 & HST-034 & 8 & 3 & 39 & 0 & 47 & $\bullet$ \\ \hline
HST-036 & 5 & HST-035 & 5 & 2 & 47 & 0 & 52 & $\bullet$ \\ \hline
HST-037 & 5 & HST-014 & 5 & 2 & 31 & 16 & 52 & \\ \hline
HST-038 & 5 & HST-007, HST-008, HST-014 & 8 & 3 & 31 & 13 & 52 & \\ \hline
\end{longtable}
}

\subsection{Diagrama de Red del Proyecto (Activity-on-Arrow)}
\begin{figure}[H]
    \centering
    \DiagramaImagen{0.95}{0.5}{imagenes/activity-on-arrow.png}
    \caption{Diagrama de red del proyecto (Activity-on-Arrow) con ruta crítica}
    \label{fig:aoa}
\end{figure}
\noindent\textbf{Ruta crítica:} HST-001 $\rightarrow$ HST-002 $\rightarrow$ HST-007 $\rightarrow$ HST-009 $\rightarrow$ HST-011 $\rightarrow$ HST-013 $\rightarrow$ HST-014 $\rightarrow$ HST-034 $\rightarrow$ HST-035 $\rightarrow$ HST-036, con una duración total de \textbf{52 días} bajo el supuesto de recursos ilimitados. Esta secuencia va desde la autenticación y el registro de almacenes, pasa por la moderación del Administrador, la búsqueda y el detalle del almacén, la reserva y el pago simulado, y finaliza en el módulo de integración de inventario (conexión externa, mapeo de campos y visualización), que concentra las tres actividades más largas del proyecto (8, 8 y 5 días).

\subsection{Planificación con Recursos Limitados}
Dado que la ruta crítica de 52 días asume disponibilidad ilimitada de personal, se elaboró además una planificación con recursos limitados, restringida a un máximo de \textbf{4 integrantes} trabajando en paralelo (el tamaño real del equipo). Una primera versión sin optimizar extiende la duración total del proyecto a \textbf{87 días}. Tras nivelar la carga de trabajo y reordenar las actividades no críticas dentro de su holgura disponible, se obtuvo un cronograma optimizado de \textbf{85 días}, que ahorra 2 días respecto a la versión sin optimizar sin exceder el límite de 4 personas ni retrasar ninguna actividad de la ruta crítica.

\begin{figure}[H]
    \centering
    \DiagramaImagen{0.95}{0.45}{imagenes/resoruce-limited-without-optimizatio.png}
    \caption{Planificación con recursos limitados sin optimizar (límite = 4 personas, duración total = 87 días)}
    \label{fig:rl-noopt}
\end{figure}

\begin{figure}[H]
    \centering
    \DiagramaImagen{0.95}{0.45}{imagenes/resource-limited-optimization.png}
    \caption{Planificación con recursos limitados optimizada (límite = 4 personas, duración total = 85 días)}
    \label{fig:rl-opt}
\end{figure}

\subsection{Diagrama de Gantt}
\begin{figure}[H]
    \centering
    \DiagramaImagen{0.95}{0.5}{imagenes/diagrama_gantt.png}
    \caption{Cronograma del proyecto (diagrama de Gantt)}
    \label{fig:gantt}
\end{figure}
\noindent\textbf{Herramienta utilizada:} Generado programáticamente (Python/Matplotlib) a partir de los datos EET/duración/ruta crítica de la Tabla~\ref{tab:cronograma_tabla}.
\clearpage
